\subsection{\texorpdfstring{例 II：\(\mathbf{B}_n\) 型与 \(\mathbf{D}_n\) 型紧李代数}{例 II：\textbackslash mathbf\{B\}\_n 型与 \textbackslash mathbf\{D\}\_n 型紧李代数}}

设 V 是有限维实向量空间，Q 是 V 上分离的正二次型。与 Q 相伴的正交群（《代数学》，第 IX 章）是子群 \(\mathbf{O}(\mathbf{Q})\) （\(\mathbf{GL}(\mathbf{V})\) 中由实希尔伯特空间 (V, Q) 的自同构组成的）；这是 \(\mathbf{GL}(\mathbf{V})\) 的李子群，其李代数是子代数 \(\mathfrak{o}(\mathbf{Q})\) （\(\mathfrak{gl}(\mathbf{V})\) 中由满足 \(x^{*} = -x\) 的 V 的自同态 x 组成的）（第 III 章，§3，第 10 节，命题 37 的推论 2），其中 \(x^{*}\) 表示 x 相对于 Q 的伴随。由于群 \(\mathbf{O}(\mathbf{Q})\) 紧，\(\mathfrak{o}(\mathbf{Q})\) 是紧实李代数。令 \(\mathbf{SO}(\mathbf{Q}) = \mathbf{O}(\mathbf{Q}) \cap \mathbf{SL}(\mathbf{V})\) ；这是 \(\mathbf{O}(\mathbf{Q})\) 的有限指数闭子群（当 dim V ≠ 0 时指数为 2），因而也以 \(\mathfrak{o}(\mathbf{Q})\) 为李代数。

当 \(V = R^{n}\) 且 Q 是通常的二次型（\(R^{n}\) 的典则基关于它标准正交）时，我们把 \(\mathbf{O}(n, \mathbf{R})\) 、\(\mathbf{SO}(n, \mathbf{R})\) 、\(\mathfrak{o}(n, \mathbf{R})\) 分别写作 \(\mathbf{O}(\mathbf{Q})\) 、\(\mathbf{SO}(\mathbf{Q})\) 、\(\mathfrak{o}(\mathbf{Q})\) 。\(\mathbf{O}(n, \mathbf{R})\)（相应地，\(\mathfrak{o}(n, \mathbf{R})\)）的元素

是矩阵 \(A \in \mathrm{M}_n(\mathbf{R})\) （满足 \(A.^t A = I_n\) ，相应地 \(A = -^t A\) ），分别称为正交的（相应地，反对称的）。

设 \(\mathrm{V}_{(\mathbf{C})}\) 是与 \(\mathrm{V}\) 相伴的复向量空间，\(\mathrm{Q}_{(\mathbf{C})}\) 是 \(\mathrm{V}_{(\mathbf{C})}\) 上与 \(\mathbf{Q}\) 相伴的二次型。把 \(\mathfrak{gl}(\mathrm{V})_{(\mathbf{C})}\) 与 \(\mathfrak{gl}(\mathrm{V}_{(\mathbf{C})})\) 等同起来；则 \(\mathfrak{o}(\mathrm{Q})_{(\mathbf{C})}\) 与 \(\mathfrak{o}(\mathrm{Q}_{(\mathbf{C})})\) 等同：这一点是显然的，因为映射 \(x \mapsto x^{*} + x\) （从 \(\mathfrak{gl}(\mathrm{V}_{(\mathbf{C})})\) 到其自身的）是 \(\mathbf{C}\) -线性的。由于 \(\mathfrak{o}(\mathrm{Q}_{(\mathbf{C})})\) 是 \(\mathrm{B}_n\) 型的（当 \(\dim V = 2n + 1\) 、\(n \geq 1\) 时），而是 \(\mathrm{D}_n\) 型的（当 \(\dim V = 2n\) 、\(n \geq 3\) 时）（第 VIII 章，§13，第 2 节与第 4 节），我们推出：

命题 5. \(B_{n}\) 型（\(n \geq 1\) ）、相应地 \(D_{n}\) 型（\(n \geq 3\) ）的每个紧单实李代数，同构于 \(\mathfrak{o}(2n + 1, \mathbf{R})\) 、相应地 \(\mathfrak{o}(2n, \mathbf{R})\) 。

